\documentclass{ximera}

\input{../../preamble.tex}

\title{Equations}

\begin{document}

\begin{abstract}
solving
\end{abstract}
\maketitle





\begin{example}

Solve $2 \sin^2(x) + 3 \sin(x) + 1 = 0$

\textbf{\textcolor{red!75!green}{explanation}} 



\[  2 \sin^2(x) + 3 \sin(x) + 1 = 0 \]

\[  (2 \sin(x) + 1) (\sin(x) + 1) = 0 \]


Either $2 \sin(x) + 1 = 0$ or $\sin(x) + 1 = 0$.


Either $-\frac{1}{2} = \sin(x) $ or $-1 = \sin(x)$.


On the interval $[0, 2\pi)$, we get three solutions.


\[ \frac{7\pi}{6}, \frac{11\pi}{6}, \frac{3\pi}{2}  \]


These can be extended with a period of $2\pi$ to get all of the solutions.






\end{example}





















\begin{center}
\textbf{\textcolor{green!50!black}{ooooo-=-=-=-ooOoo-=-=-=-ooooo}} \\

more examples can be found by following this link\\ \link[More Examples of Trigonometric Functions]{https://ximera.osu.edu/csccmathematics/precalculus/precalculus/trigForms/examples/exampleList}

\end{center}







\end{document}
